% 第3回　先行研究の探索（記入例・模範解答）
\session{3}{1}{10月14日（水）}{先行研究の探索}{AI検索の活用}{}

\goals{
  \item AIを使って関連分野と検索キーワード（日本語と英語）を広げられる
  \item 学術データベースで文献を探し、実在と内容を確認できる
  \item 文献リストに必要な情報（著者、発行年、題名、掲載誌または出版社、巻号と頁、DOIまたはURL）を揃えられる
}

\afillin{自分のリサーチ・クエスチョン}{名古屋都市圏の20代にとって、オンラインでの情報収集は、ディーラーへの来店のしかたをどう変えたのか}

\work[60mm]{ワーク1}{探す場所の見当をつける}
\begin{promptbox}[探す場所の見当をつける]
次のリサーチ・クエスチョンに関連する学術分野と、文献検索に使うキーワードを日本語と英語でそれぞれ10個ずつ挙げてください。文献名や著者名は挙げないでください。問い：（あなたの問い）
\end{promptbox}
\afillin{関連する学術分野}{マーケティング（消費者行動論）、流通論、交通経済学、情報行動論}
\begin{wstabh}{colspec={X[l,m] X[l,m] X[l,m] X[l,m]},row{2-Z}={ht=8mm}}
  日本語キーワード & 英語キーワード & 日本語キーワード & 英語キーワード \\
  \af{若者の車離れ} & \af{young adults car ownership} & \af{消費者の情報探索} & \af{consumer information search} \\
  \af{自動車ディーラー} & \af{car dealership} & \af{購買意思決定} & \af{purchase decision} \\
  \af{来店行動} & \af{store visit behavior} & \af{オムニチャネル} & \af{omnichannel retailing} \\
  \af{オンライン購買} & \af{online car buying} & \af{口コミ} & \af{online reviews} \\
  \af{カーシェアリング} & \af{car sharing} & \af{都市交通} & \af{urban mobility} \\
\end{wstabh}

\work[50mm]{ワーク2}{データベースで検索する}
\instr{CiNii Research、Google Scholar、J-STAGE、図書館の蔵書検索などを使う。検索の記録を残す。}
\begin{wstabh}{colspec={Q[l,wd=28mm,m] X[2,l,m] Q[c,wd=14mm,m] X[3,l,m]},row{2-Z}={ht=10mm}}
  データベース & 検索語 & 件数 & 関係の深そうな文献 \\
  \af{CiNii Research} & \af{若者 車離れ} & \af{（記録）} & \af{若者の自動車保有に関する論文。題名と要旨をメモ} \\
  \af{CiNii Research} & \af{自動車 販売店 来店} & \af{（記録）} & \af{販売店の顧客対応に関する論文} \\
  \af{J-STAGE} & \af{消費者 情報探索 インターネット} & \af{（記録）} & \af{購買前の情報探索に関する論文} \\
  \af{Google Scholar} & \af{online information search car purchase} & \af{（記録）} & \af{海外の自動車購入の情報探索に関する論文} \\
  \af{官公庁サイト} & \af{乗用車市場動向調査／通信利用動向調査} & \af{—} & \af{公的調査の報告書（下の文献リスト）} \\
\end{wstabh}

\work[60mm]{ワーク3}{AIが挙げた文献の実在を確かめる}
\instr{AIに「文献を挙げて」と頼み、挙がったものをデータベースで確認する。}
\begin{wstabh}{colspec={X[3,l,m] Q[c,wd=13mm,m] Q[c,wd=13mm,m] X[2,l,m]},row{2-Z}={ht=10mm}}
  AIが挙げた文献 & 実在\newline ○× & 内容\newline 一致 & 確認先・違っていた点 \\
  \af{若者の車離れの要因を分析したとされる論文（著者・学会誌名・年つき）} & \ans{×} & \ans{—} & \af{CiNii・J-STAGE とも該当なし} \\
  \af{自動車販売のデジタル化を扱ったとされる書籍} & \ans{×} & \ans{—} & \af{蔵書検索・CiNii Books で該当なし} \\
  \af{業界団体の乗用車市場の調査報告} & \ans{○} & \ans{△} & \af{調査は実在。AIが示した数値は報告書に見当たらない} \\
  \af{総務省の情報通信白書} & \ans{○} & \ans{○} & \af{総務省サイトで確認。版と章を特定した} \\
  \af{海外の学術誌に載ったとされる英語論文} & \ans{○} & \ans{×} & \af{論文は実在するが、テーマは自動車ではなく家電} \\
\end{wstabh}
{\small 集計：挙がった文献 \underline{\makebox[10mm]{\ans{5}}} 件のうち、実在しない \underline{\makebox[10mm]{\ans{2}}} 件、内容が違う \underline{\makebox[10mm]{\ans{2}}} 件}\par

\work[80mm]{まとめ}{文献リスト（実在を確認した5件）}
\begin{wstabh}{colspec={Q[c,wd=6mm,m] X[2,l,m] Q[c,wd=10mm,m] X[3,l,m] X[3,l,m] X[2,l,m]},row{2-Z}={ht=15mm}}
  & 著者 & 年 & 題名 & 掲載誌・出版社、巻号・頁 & DOI・URL \\
  1 & \af{日本自動車工業会} & \af{※} & \af{乗用車市場動向調査} & \af{日本自動車工業会（隔年の調査報告書）} & \af{jama.or.jp} \\
  2 & \af{総務省} & \af{※} & \af{情報通信白書} & \af{総務省（年次の白書）} & \af{soumu.go.jp} \\
  3 & \af{総務省} & \af{※} & \af{通信利用動向調査} & \af{総務省（年次の調査報告）} & \af{soumu.go.jp} \\
  4 & \af{経済産業省} & \af{※} & \af{電子商取引に関する市場調査} & \af{経済産業省（年次の調査報告書）} & \af{meti.go.jp} \\
  5 & \af{国土交通省} & \af{※} & \af{全国都市交通特性調査} & \af{国土交通省（定期の調査）} & \af{mlit.go.jp} \\
\end{wstabh}
{\footnotesize\color{ans}※ 記入例では年次・版・URLの詳細を省いた。学生には、参照した版の発行年と、該当ページのURLまで書かせる。}\par

\begin{nexttime}
  \item 実在を確認した文献5件のリストを完成させる。うち1件は本文を入手し、第4回に持ってくる
  \item AIが挙げた文献のうち、実在しなかった・内容が違ったものの数と例をメモしておく（上のワーク3）
\end{nexttime}
\aimemoA{関連分野・キーワードの洗い出し、文献の実在確認の練習}
{「関連する学術分野と、検索キーワードを日本語と英語で10個ずつ。文献名は挙げないで」}
{キーワードは、データベースで実際にヒットしたものだけを使った。AIが挙げた文献5件は、確認できたものだけを残し、リストには自分で探した文献を使った。}
{CiNii Research、J-STAGE、大学図書館の蔵書検索、各省庁と業界団体の公式サイト}

\begin{point}
  \item 文献リストの必須項目（著者、発行年、題名、掲載誌または出版社、巻号と頁、DOIまたはURL）が揃っているかを見る。年やURLの欠けが最も多い。
  \item 記入例の文献は、実在が確実な公的統計・白書に限っている。学生には、学術論文（CiNii・J-STAGE で見つけたもの）を少なくとも2件含めるよう指導する。この冊子では架空の論文を「記入例」として載せることを避けた。
  \item ワーク3の集計は、AIの出力によって大きく変わる。数の多少ではなく、「実在しても内容が違う」ケースに気づけたかを評価する。
  \item 検索の件数は検索日によって変わるため、記入例では空欄にした。検索日と件数を記録する習慣をつけさせる。
\end{point}
